\chapter{投资建议}

\begin{houseviewbox}[投资建议总纲]
基于 2026 年 7 月 3 日收盘价统一快照，本报告覆盖的 56 只可发布目标价/公允价值标的的市值加权综合目标空间为 \textbf{-20.6\%}，较 6 月末快照的 -32.1\% 显著改善。9 只标的已具备正收益空间，其中 \textbf{明阳电气（+72.7\%）、中兴通讯（+56.6\%）、星网锐捷（+52.2\%）} 证据链可审计且上行空间最大。投资动作以"核心关注、回调验证、市场支撑观察、高估值风险"四级标签为主，不把主题热度当作买入理由。当前阶段更适合事件验证和回撤布局，而非无差别追高。
\end{houseviewbox}

\section{投资评级与目标价全景}

\begin{exhibitbox}[56 只标的投资评级全景表（按综合目标空间排序）]
\small
\begin{longtable}{L{1.5cm}L{2.0cm}R{1.3cm}R{1.3cm}R{1.3cm}R{1.5cm}R{1.3cm}L{1.8cm}}
\toprule
\textbf{代码} & \textbf{公司} & \textbf{现价} & \textbf{2026E EPS} & \textbf{基础锚} & \textbf{综合目标} & \textbf{空间} & \textbf{动作}\\
\midrule
\endhead
\multicolumn{8}{l}{\textit{核心关注（正收益空间，证据可审计）}}\\
301291 & 明阳电气 & 36.19 & 3.30 & 82.5 & 62.5 & +72.7\% & 核心关注\\
000063 & 中兴通讯 & 36.38 & 1.94 & 73.7 & 57.0 & +56.6\% & 核心关注\\
002396 & 星网锐捷 & 23.40 & 1.21 & 46.0 & 35.6 & +52.2\% & 核心关注\\
\midrule
\multicolumn{8}{l}{\textit{回调验证（等待更好安全边际）}}\\
002913 & 奥士康 & 57.15 & 2.18 & 82.8 & 67.8 & +18.7\% & 回调验证\\
600089 & 特变电工 & 22.01 & 1.21 & 30.2 & 25.5 & +15.9\% & 回调验证\\
600941 & 中国移动 & 87.48 & 7.24 & 115.8 & 96.7 & +10.6\% & 回调验证\\
600050 & 中国联通 & 4.07 & 0.34 & 5.4 & 4.5 & +10.6\% & 回调验证\\
601728 & 中国电信 & 5.43 & 0.42 & 6.7 & 5.7 & +4.9\% & 回调验证\\
000530 & 冰山冷热 & 5.36 & 0.25 & 6.2 & 5.4 & +0.3\% & 回调验证\\
\midrule
\multicolumn{8}{l}{\textit{市场支撑观察（情绪流动性支撑，等待财务确认）}}\\
601138 & 工业富联 & 64.72 & 2.21 & 75.1 & 64.1 & -0.9\% & 市场支撑观察\\
000938 & 紫光股份 & 30.28 & 1.20 & 29.9 & 28.8 & -4.9\% & 市场支撑观察\\
002922 & 伊戈尔 & 29.65 & 1.15 & 28.7 & 27.7 & -6.7\% & 市场支撑观察\\
300502 & 新易盛 & 526.00 & 12.17 & 487.0 & 479.8 & -8.8\% & 市场支撑观察\\
300308 & 中际旭创 & 1116.00 & 22.52 & 1013.5 & 1001.1 & -10.3\% & 市场支撑观察\\
002158 & 汉钟精机 & 35.97 & 1.34 & 33.5 & 31.3 & -12.9\% & 市场支撑观察\\
603228 & 景旺电子 & 71.97 & 2.18 & 69.8 & 62.5 & -13.2\% & 市场支撑观察\\
002518 & 科士达 & 47.80 & 1.59 & 39.8 & 40.0 & -16.4\% & 市场支撑观察\\
600845 & 宝信软件 & 18.25 & 1.03 & 16.5 & 15.2 & -16.8\% & 市场支撑观察\\
\midrule
\multicolumn{8}{l}{\textit{高估值风险（价格透支可验证增长）}}\\
000977 & 浪潮信息 & 66.35 & 1.65 & 46.1 & 49.6 & -25.3\% & 高估值风险\\
300476 & 胜宏科技 & 308.07 & 5.92 & 189.4 & 229.1 & -25.6\% & 高估值风险\\
002335 & 科华数据 & 37.95 & 0.93 & 23.2 & 27.5 & -27.6\% & 高估值风险\\
688123 & 聚辰股份 & 204.00 & 3.49 & 132.6 & 144.8 & -29.0\% & 高估值风险\\
002463 & 沪电股份 & 135.35 & 2.58 & 90.4 & 95.4 & -29.5\% & 高估值风险\\
601179 & 中国西电 & 14.29 & 0.40 & 9.9 & 10.0 & -29.9\% & 高估值风险\\
300990 & 同飞股份 & 100.10 & 2.76 & 69.0 & 69.4 & -30.7\% & 高估值风险\\
300394 & 天孚通信 & 250.10 & 2.75 & 151.4 & 169.7 & -32.1\% & 高估值风险\\
688256 & 寒武纪 & 1353.00 & 16.10 & 611.8 & 907.9 & -32.9\% & 高估值风险\\
301165 & 锐捷网络 & 94.89 & 1.44 & 54.7 & 62.8 & -33.8\% & 高估值风险\\
002837 & 英维克 & 71.43 & 0.64 & 22.3 & 45.8 & -35.9\% & 高估值风险\\
300442 & 润泽科技 & 85.44 & 1.57 & 47.0 & 54.2 & -36.5\% & 高估值风险\\
300738 & 奥飞数据 & 20.36 & 0.39 & 9.4 & 12.8 & -37.2\% & 高估值风险\\
000988 & 华工科技 & 153.95 & 1.99 & 75.6 & 93.6 & -39.2\% & 高估值风险\\
000021 & 深科技 & 55.91 & 0.64 & 24.2 & 33.0 & -41.0\% & 高估值风险\\
603083 & 剑桥科技 & 209.37 & 2.52 & 95.8 & 122.3 & -41.6\% & 高估值风险\\
688676 & 金盘科技 & 79.54 & 1.14 & 31.8 & 43.3 & -45.5\% & 高估值风险\\
688183 & 生益电子 & 122.89 & 1.46 & 46.7 & 65.6 & -46.6\% & 高估值风险\\
600183 & 生益科技 & 160.78 & 2.00 & 60.0 & 84.7 & -47.3\% & 高估值风险\\
002281 & 光迅科技 & 217.33 & 1.90 & 72.2 & 114.5 & -47.3\% & 高估值风险\\
300499 & 高澜股份 & 36.60 & 0.45 & 11.2 & 19.2 & -47.6\% & 高估值风险\\
688008 & 澜起科技 & 266.80 & 1.92 & 73.0 & 138.8 & -48.0\% & 高估值风险\\
688041 & 海光信息 & 325.40 & 2.01 & 76.4 & 168.2 & -48.3\% & 高估值风险\\
603986 & 兆易创新 & 677.77 & 4.45 & 169.1 & 337.6 & -50.2\% & 高估值风险\\
002364 & 中恒电气 & 52.90 & 0.65 & 16.2 & 25.7 & -51.3\% & 高估值风险\\
002916 & 深南电路 & 454.48 & 5.33 & 149.3 & 219.0 & -51.8\% & 高估值风险\\
603186 & 华正新材 & 199.99 & 1.77 & 56.5 & 94.1 & -52.9\% & 高估值风险\\
603019 & 中科曙光 & 92.68 & 0.71 & 29.8 & 42.5 & -54.1\% & 高估值风险\\
301018 & 申菱环境 & 117.78 & 0.91 & 27.2 & 52.2 & -55.7\% & 高估值风险\\
300223 & 北京君正 & 259.56 & 1.36 & 51.7 & 113.5 & -56.3\% & 高估值风险\\
688519 & 南亚新材 & 371.22 & 2.08 & 66.6 & 149.8 & -59.6\% & 高估值风险\\
688521 & 芯原股份 & 303.00 & 0.17 & 6.5 & 120.5 & -60.2\% & 高估值风险\\
301205 & 联特科技 & 297.99 & 1.28 & 48.6 & 117.1 & -60.7\% & 高估值风险\\
300249 & 依米康 & 14.82 & 0.13 & 3.2 & 5.7 & -61.2\% & 高估值风险\\
300383 & 光环新网 & 13.00 & 0.05 & 0.8 & 4.5 & -65.0\% & 高估值风险\\
603881 & 数据港 & 25.21 & 0.23 & 3.7 & 8.6 & -66.0\% & 高估值风险\\
688702 & 盛科通信 & 323.64 & 0.11 & 4.2 & 95.6 & -70.5\% & 高估值风险\\
688795 & 摩尔线程 & 643.81 & -2.13 & 182.2 & 182.2 & -71.7\% & 高估值风险\\
\bottomrule
\end{longtable}
\sourcenote{Sina 2026-07-03 收盘价；akshare 财务摘要；AStock 估值模型 combined\_target\_valuation\_model\_20260703。空间=综合目标价/现价-1。56 只标的市值加权平均空间 -24.1\%。}
\end{exhibitbox}

\begin{riskbox}[动作标签定义]
\textbf{核心关注}：当前目标空间为正（>5\%），证据链可审计（直接收入证据 + 财务交付 + 客户链验证），优先纳入组合跟踪池。\\
\textbf{回调验证}：公司质量较好但当前价格需要更好安全边际，等待回调至目标价附近或中报业绩上修后再评估。\\
\textbf{市场支撑观察}：现价主要由情绪和流动性支撑，基本面证据需等待中报或订单确认才能支撑当前估值。\\
\textbf{高估值风险}：当前价格已显著透支模型可验证增长，若中报不能同时兑现收入增速、毛利率和现金流，当前价格隐含的增长久期会迅速收缩。
\end{riskbox}

\section{核心关注标的深度}

\begin{exhibitbox}[核心关注标的——明阳电气（301291）]
\textbf{投资逻辑}：公司是新能源输配电翘楚，海外业务和 AIDC 供配电带来新增长曲线。2025 年上半年业绩高增，AIDC 变压器和高低压柜业务在运营商和互联网客户中取得突破。当前估值仅反映传统业务价值，AIDC 业务贡献期权价值未被定价。\\
\textbf{证据链}：海外营收 +68\%，新签订单 +281\%，高压变压器订单有望超预期。AIDC 相关订单已在运营商集采中中标。\\
\textbf{目标价}：62.5 元（基础锚 82.5 元，综合目标 62.5 元），上行空间 +72.7\%。\\
\textbf{关键催化剂}：运营商 AIDC 供配电集采中标、海外大单交付、中报业绩确认。\\
\textbf{下行风险}：新能源业务承压、原材料价格上涨、AIDC 项目交付延迟。
\end{exhibitbox}

\begin{exhibitbox}[核心关注标的——中兴通讯（000063）]
\textbf{投资逻辑}：AI 交换机和服务器业务在运营商和互联网客户中加速渗透，公司在 800G 交换机和 DCI 设备上的技术积累开始转化为收入。当前估值仅反映通信设备业务价值，AI 网络设备的增长期权被严重低估。\\
\textbf{证据链}：2026Q1 营收 +6.1\%，研发投入持续加大。AI 交换机产品在运营商集采中份额提升，服务器业务受益国产算力替代。\\
\textbf{目标价}：57.0 元（基础锚 73.7 元，综合目标 57.0 元），上行空间 +56.6\%。\\
\textbf{关键催化剂}：运营商 AI 交换机集采结果、国产算力服务器订单、中报研发投入转化效率。\\
\textbf{下行风险}：研发投入拖累短期利润、海外地缘风险、AI 交换机价格竞争。
\end{exhibitbox}

\begin{exhibitbox}[核心关注标的——星网锐捷（002396）]
\textbf{投资逻辑}：数据中心交换机驱动利润高增，公司在 AI 数据中心网络设备和 KVM 领域的市场份额持续提升。探索跨境支付应用带来额外期权价值。当前估值显著低于同业网络设备公司。\\
\textbf{证据链}：2025 年中报预告点评确认数据中心交换机驱动利润高增。Q1 净利润同比高增，ICT 基建和 AI 应用布局加强。\\
\textbf{目标价}：35.6 元（基础锚 46.0 元，综合目标 35.6 元），上行空间 +52.2\%。\\
\textbf{关键催化剂}：数据中心交换机订单、KVM 在 AI 数据中心渗透率提升、跨境支付业务进展。\\
\textbf{下行风险}：交换机价格竞争、客户集中度风险、新业务拓展不及预期。
\end{exhibitbox}

\section{分链条投资策略}

\begin{exhibitbox}[八大链条投资策略矩阵]
\begin{tabularx}{\textwidth}{L{2.2cm}L{2.5cm}X X}
\toprule
\textbf{链条} & \textbf{首选标的} & \textbf{投资逻辑} & \textbf{风险点}\\
\midrule
算力芯片/存储 & 澜起科技（观察） & 内存接口芯片直接受益 DDR5 渗透加速，Q1 净利 +61\%，毛利率 69.8\% 行业最高 & 海光信息、寒武纪估值过高（PE 167x/87x），基本面锚贡献不足 40\%\\
服务器/交换机 & 工业富联、中兴通讯 & AI 服务器 ODM 龙头证据链完整（Q1 GPU 机柜 +380\%）；中兴通讯 AI 交换机被低估 & 浪潮信息 Q1 营收 -24\%，交付节奏待观察\\
网络/光通信 & 中际旭创、新易盛 & 800G/1.6T 出货量直接证据最强，Q1 业绩超预期（旭创净利 +262\%） & ASP 下滑风险（VP-02），拥挤交易踩踏风险（VP-05）\\
PCB/材料 & 胜宏科技（谨慎）、景旺电子 & 高端 AI PCB 需求真实，但胜宏科技估值已透支；景旺电子估值相对合理 & 扩产和良率压力侵蚀毛利，二线 PCB 产能利用率承压\\
供配电 & 明阳电气、特变电工 & 明阳电气 AIDC 供配电突破 + 海外高增；特变电工估值安全边际好 & 电力瓶颈（VP-03）限制 IDC 建设速度，供配电需求可能延后\\
液冷/温控 & 英维克（谨慎）、高澜股份 & 政策驱动液冷渗透率提升逻辑真实，但当前估值隐含 30\%+ 渗透率（实际可能仅 10-15\%） & 液冷渗透延迟（VP-06），运维复杂度和标准缺失阻碍大规模部署\\
IDC 运营 & 三大运营商 & 运营商电力指标资源更优，防御属性强，当前已具正收益空间 & 润泽科技、奥飞数据规划 MW 远大于已获批 MW，电力指标审批延迟\\
下游需求 & 宝信软件（观察） & 制造 AI 和工业互联网需求稳健，但 AIDC 直接收入占比低 & 需求幻觉（VP-01）：CSP capex 兑现率可能仅 60-70\%\\
\bottomrule
\end{tabularx}
\sourcenote{analysis/chain\_business\_research.md；analysis/valuation\_model.md；improvements/variant\_perception\_7\_arguments.md。}
\end{exhibitbox}

\section{组合构建建议}

\begin{exhibitbox}[AIDC 产业链组合构建框架]
\textbf{组合目标}：在控制下行风险的前提下，捕捉 AIDC 产业链盈利周期从"主题预期"进入"财报验证"阶段的结构性机会。\\
\textbf{基准权重}（建议初始配置 10-15\% 组合权重，视风险偏好调整）：
\begin{itemize}
\item \textbf{核心仓位（40\%）}——明阳电气、中兴通讯、星网锐捷：正收益空间 + 证据可审计，作为组合底仓
\item \textbf{质量仓位（30\%）}——工业富联、中际旭创、新易盛：证据链最完整但估值偏高，等待更好入场点
\item \textbf{防御仓位（20\%）}——中国移动、中国联通、中国电信：电力指标优势 + 防御属性，对冲板块回调
\item \textbf{期权仓位（10\%）}——寒武纪、兆易创新：业绩上修弹性大但估值风险高，小仓位参与中报催化
\end{itemize}
\end{exhibitbox}

\begin{riskbox}[组合再平衡规则]
\textbf{月度再平衡}：根据月度价格变动和证据更新调整权重，单一标的权重不超过组合 8\%。\\
\textbf{事件驱动调整}：中报业绩预告超预期的标的可临时提权 2-3 个百分点；业绩不及预期的立即降权至观察池。\\
\textbf{止损规则}：核心仓位设置 15\% 止损线，质量仓位设置 20\% 止损线，期权仓位设置 25\% 止损线。\\
\textbf{拥挤度信号}：当主动基金 AIDC 配置比例 >12\% 或融资余额/流通市值 >3.5\% 时，整体减仓 1/3。
\end{riskbox}

\section{催化剂与监测框架}

\begin{exhibitbox}[下一季度关键催化剂与上下行触发]
\begin{tabularx}{\textwidth}{L{2.8cm}X X}
\toprule
\textbf{时间窗口} & \textbf{上行触发} & \textbf{下行触发}\\
\midrule
7 月中旬-8 月上旬 & 中报业绩预告超预期：寒武纪 +20-30\%、兆易创新 +30-50\%、中际旭创 +10-15\% & 中报不及预期：光模块 ASP 下滑、AI 服务器增收不增利、PCB 产能利用率下降\\
7 月下旬 & 北美 CSP Q2 财报及全年 capex 指引上修 & 北美 CSP capex 指引下修 15\%+（VP-01 核心证伪条件）\\
8 月 & 运营商 AI 服务器/液冷集采结果公布 & 运营商集采液冷比例 <20\%（VP-06 证伪条件）\\
9 月 & 1.6T 光模块客户验证进展、国产 AI 芯片大模型部署 & 800G ASP 环比下滑 >8\%（VP-02 证伪条件）\\
持续监测 & 能耗指标审批加速、绿电 PPA 签约 & 能耗审批暂停或收紧（VP-03 强化条件）\\
\bottomrule
\end{tabularx}
\sourcenote{analysis/chain\_earnings\_bridge.md；improvements/variant\_perception\_7\_arguments.md；analysis/2026\_h1\_earnings\_pre\_announcement\_update.md。}
\end{exhibitbox}

\begin{exhibitbox}[高频监测指标体系]
\begin{tabularx}{\textwidth}{L{2.5cm}L{2.0cm}L{2.0cm}X}
\toprule
\textbf{指标} & \textbf{频率} & \textbf{当前值} & \textbf{预警阈值}\\
\midrule
北美 CSP capex 增速 & 季度 & +40-60\% YoY & 连续两季 <20\%\\
800G ASP & 月度 & \$1800-2200 & 环比下滑 >5\%\\
液冷服务器招标占比 & 月度 & 18-22\% & <15\%\\
AIDC 板块融资余额/流通市值 & 日度 & 2.8-3.2\% & >3.5\%\\
主动基金 AIDC 配置比例 & 季度 & 8-10\% & >12\%\\
北向资金光模块持仓 & 日度 & 8-12\% & 连续 5 日净流出\\
IDC 能耗指标审批量 & 月度 & 约 750MW/年 & 环比下滑 >30\%\\
\bottomrule
\end{tabularx}
\sourcenote{improvements/variant\_perception\_7\_arguments.md；analysis/risk\_framework.md。}
\end{exhibitbox}

\section{反方观点与对冲策略}

\begin{houseviewbox}[七大反方论点与置信度]
当前 AIDC 共识叙事在 7 个维度上存在可被证伪的脆弱性。按置信度排序：
\begin{enumerate}
\item \textbf{估值锚失效（85\%置信度）}——盛科通信（3372x PE）、芯原股份（1904x PE）等标的基本面锚贡献不足 5\%，目标价无基本面意义
\item \textbf{国产替代伪命题（80\%置信度）}——A 股在 GPU/HBM/DSP 等核心价值池缺乏高纯度标的，利润池被海外拿走
\item \textbf{电力天花板（75\%置信度）}——规划项目电力需求远超电网可供给增量，约 2/3 规划项目可能无法获批
\item \textbf{需求幻觉（70\%置信度）}——CSP capex 兑现率历史仅 60-70\%，当前模型假设 100\% 兑现
\item \textbf{液冷渗透延迟（70\%置信度）}——2026 年液冷渗透率可能仍低于 15\%，远低于市场预期的 30\%+
\item \textbf{ASP 崩塌（65\%置信度）}——800G 价格下滑速度可能超预期，二线厂商产能投放 + CSP 压价
\item \textbf{拥挤交易踩踏（60\%置信度）}——光模块三巨头合计 3.5 万亿市值，机构持仓集中度历史新高
\end{enumerate}
\end{houseviewbox}

\begin{exhibitbox}[综合失败场景与对冲策略]
\textbf{触发事件}（2026 年 8-10 月概率 25-30\%）：北美 CSP Q3 财报电话会下调全年 capex 指引 15-20\%，同时 800G ASP 环比下滑 8-10\%。\\
\textbf{板块影响}：光模块 -40\~50\%、AI 服务器 -30\~40\%、PCB -25\~35\%、液冷 -35\~45\%、IDC 运营 -20\~30\%、国产算力 -45\~55\%。\\
\textbf{对冲策略}：
\begin{enumerate}
\item 减配光模块至标配以下 3 个百分点
\item 增配三大运营商（电力指标优势 + 防御属性）至组合 5\%+
\item 买入沪深 300 看跌期权（行权价 -10\%，期限 6 个月）
\item 对高 PE 标的（盛科通信、芯原股份、寒武纪）设置 15\% 止损
\item 对基本面锚贡献 <10\% 的标的，目标价仅作"市场支撑观察"参考，不作投资建议依据
\end{enumerate}
\sourcenote{improvements/variant\_perception\_7\_arguments.md。}
\end{exhibitbox}

\begin{riskbox}[ESG/HRF 筛查与排除清单]
经 ESG/HRF 筛查，以下标的因 BIS 实体清单或其他重大治理风险需从核心池排除：
\begin{itemize}
\item \textbf{中科曙光（603019）}——BIS 实体清单，从核心候选池排除
\item \textbf{海光信息（688041）}——BIS 实体清单，从核心候选池排除
\end{itemize}
3 只标的基本面锚定贡献不足 15\%，标注为"市场情绪主导型"：盛科通信（688702）、芯原股份（688521）、光环新网（300383）。
\end{riskbox}

\section{观察名单与升级路径}

\begin{exhibitbox}[观察名单公司——PS/SOTP 替代估值与升级触发]
\begin{tabularx}{\textwidth}{L{1.5cm}L{1.8cm}R{1.2cm}R{1.2cm}R{1.2cm}R{1.5cm}L{2.5cm}}
\toprule
\textbf{代码} & \textbf{公司} & \textbf{现价} & \textbf{SOTP} & \textbf{期权上沿} & \textbf{差距评分} & \textbf{最可能升级时间}\\
\midrule
300936 & 中英科技 & 105.6 & 40.6 & 58.6 & 2.6/10 & 2026Q4-2027Q1\\
603912 & 佳力图 & 7.75 & 5.4 & 8.2 & 3.8/10 & 2026Q3-Q4\\
002334 & 英威腾 & 7.15 & 8.7 & 10.9 & 4.4/10 & 2026Q2-Q3\\
\bottomrule
\end{tabularx}
\vspace{0.5em}
\textbf{升级条件}（满足任一即可从观察名单升级为正式目标价模型）：
\begin{enumerate}
\item 单季度毛利率环比改善 >2pct 且 AIDC 收入同比 >50\%
\item 获得单笔 >1 亿元 AIDC 相关订单
\item 券商首次覆盖并给出明确目标价
\end{enumerate}
\sourcenote{improvements/watchlist\_valuation\_triggers.md。}
\end{exhibitbox}

\begin{houseviewbox}[最终建议总结]
\textbf{一句话结论}：AIDC 产业链盈利周期真实且在加速，但 2026 年 7 月价格水平下，覆盖组合加权平均目标空间为 -20.6\%，多数核心标的现价已反映乐观预期。\\
\textbf{优先动作}：核心关注明阳电气、中兴通讯、星网锐捷（正收益空间 + 证据可审计）；回调关注工业富联、三大运营商（质量好但需更好入场点）；回避高估值情绪驱动标的（盛科通信、芯原股份等基本面锚贡献 <10\%）。\\
\textbf{关键验证窗口}：7 月中旬-8 月上旬中报业绩预告，重点验证光模块 ASP、AI 服务器毛利率、PCB 高端产品占比。\\
\textbf{组合纪律}：初始配置 10-15\% 权重，单一标的不超过 8\%，设置 15-25\% 分层止损，拥挤度超标时减仓 1/3。
\end{houseviewbox}